\subsection{交换群正则表示的子表示}

在本节中，G 是一个局部紧交换群，\(\mu\) 是 G 上的一个哈尔测度。用 \(\hat{G}\) 表示 G 的对偶群（II, 第 201 页, 定义 2），用 \(\hat{\mu}\) 表示 \(\mu\) 在 \(\hat{G}\) 上的对偶哈尔测度（II, 第 214 页, 定义 4）。可测性概念总是相对于 \(\mu\) 与 \(\hat{\mu}\) 而言。

我们打算确定左正则表示 \(\gamma_{G}\)（G 在 \(L^{2}(G,\mu)\) 中的）的所有子表示。由于 G 是交换的，我们还有 \(\delta_{\mathrm{G}}(g)=\gamma_{\mathrm{G}}(g^{-1})\)，故这些子表示也是右正则表示的子表示。

对 \(\widehat{\mathbf{G}}\) 的每个可测子集 M，用 \(\mathrm{E_M}\) 表示满足下列条件的 \(f \in \mathrm{L}^2 (\mathrm{G},\mu)\) 组成的集合：傅里叶变换 \(\mathcal{F}_{\mathrm{G}}(f)\) 在 M 上几乎处处为零（参见 II, 第 210 页, 第 3 目）。它是连续线性映射 \(f \mapsto \varphi_{\mathrm{M}}\mathcal{F}_{\mathrm{G}}(f)\)（从 \(\mathrm{L}^2 (\mathrm{G},\mu)\) 到 \(\mathrm{L}^2 (\widehat{\mathrm{G}},\widehat{\mu})\) 中）的核，其中 \(\varphi_{\mathrm{M}}\) 表示 M 的特征函数（参见 II, 第 215 页, 定理 1），因而是 \(\mathrm{L}^2 (\mathrm{G},\mu)\) 的一个闭子空间。

若 \(\widehat{G}\) 的可测子集 M 与 N 使得 \((\mathrm{M} \cup \mathrm{N}) - (\mathrm{M} \cap \mathrm{N})\) 是局部零测的，则称它们相差一个局部零测集意义下相等。等价地，可测子集 M 与 N 相差一个局部零测集意义下相等，当且仅当 M 与 N 的特征函数在 \(\mathrm{L}^{\infty}(\widehat{\mathrm{G}}, \widehat{\mu})\) 中相等。

命题 10. --- a) 设 M 是 \(\widehat{G}\) 的一个可测子集。空间 \(\mathrm{E}_{\mathrm{M}}\) 是表示 \(\gamma_{\mathrm{G}}\) 的一个子表示；

\begin{enumerate}
\def\labelenumi{\alph{enumi})}
\setcounter{enumi}{1}
\item
  设 M 与 N 是 \(\widehat{\mathbf{G}}\) 的可测子集。我们有 \(\mathrm{E_M} = \mathrm{E_N}\) 当且仅当 M 与 N 相差一个局部零测集意义下相等；
\item
  \(\gamma_{\mathrm{G}}\) 的任何子表示都形如 \(\mathrm{E}_{\mathrm{M}}\)，其中 M 是 \(\widehat{\mathrm{G}}\) 的一个可测子集。
\end{enumerate}

用 \(\eta\) 表示 G 到 \(\widehat{\mathbf{G}}\) 的典范映射（参见 II, 第 216 页, 注 1）。由于 \(\mathrm{E_M}\) 是闭的，断言 \(a\) ) 由以下公式得出

\[
\boldsymbol {\gamma} _ {\mathrm{G}} (x) (f) = \varepsilon_ {x} * f, \quad \mathcal {F} _ {\mathrm{G}} (\varepsilon_ {x} * f) = \eta (x) \mathcal {F} _ {\mathrm{G}} (f)
\]

其中 \(x\in \mathrm{G}\)，\(f\in \mathrm{L}^2 (\mathrm{G},\mu)\)。

设 M 与 N 是相差一个局部零测集意义下相等的可测子集。由于此时 \(\varphi_{M}\) 等于 \(\varphi_{N}\)（在 \(L^{\infty}(\widehat{G},\widehat{\mu})\) 中），这一条件蕴含 \(E_{M}=E_{N}\)。

反之，设 M 与 N 不是相差一个局部零测集意义下相等的。那么在交换 M 与 N 的角色之后，存在紧子集 K 使得集合 \(L = K \cap (M - (M \cap N))\) 不是零测的。设 \(\varphi_{\mathrm{L}} \in \mathrm{L}^{2}(\widehat{\mathrm{G}}, \widehat{\mu})\) 是 L 的特征函数所在的类，并令 \(f = \overline{\mathcal{F}}_{\widehat{\mathrm{G}}}(\varphi_{\mathrm{L}}) \in \mathrm{L}^{2}(\mathrm{G}, \mu)\)；我们有 \(\mathcal{F}_{\mathrm{G}}(f) = \varphi_{\mathrm{L}}\)（II, 第 220 页, 定理 2 的推论）。于是 f 属于 \(E_{N}\)，因为 \(L \cap N\) 是空集；但不属于 \(E_{M}\)，因为 \(\varphi_{\mathrm{M}}\mathcal{F}(f) = \varphi_{\mathrm{M}}\varphi_{\mathrm{L}} = \varphi_{\mathrm{L}}\)。因此 \(E_{M} \neq E_{N}\)，这就证明了 b)。

现在设 E 是 \(\gamma_{\mathrm{G}}\) 的一个子表示。设 \(p_{\mathrm{E}}\) 是 \(\mathrm{L}^2 (\mathrm{G},\mu)\) 中以 E 为像的正交投影算子，并令 \(q_{\mathrm{E}} = \mathcal{F}_{\mathrm{G}}\circ p_{\mathrm{E}}\circ \mathcal{F}_{\mathrm{G}}^{-1}\)。投影算子 \(p_{\mathrm{E}}\) 属于 \(\mathrm{Hom}_{\mathrm{G}}(\gamma_{\mathrm{G}},\gamma_{\mathrm{G}})\)（V, 第 383 页, 命题 4），故它与 \(\gamma_{\mathrm{G}}(f)\) 交换，对所有 \(f\in \mathrm{L}^{1}(\mathrm{G},\mu)\)（参见 V, 第 401 页）。这意味着它与自同态 \(\varphi \mapsto f*\varphi\) 交换，对所有 \(f\in \mathrm{L}^{1}(\mathrm{G},\mu)\)。于是由 V 第 407 页的引理 4，自同态 \(q_{\mathrm{E}}\)（\(\mathrm{L}^2 (\widehat{\mathrm{G}},\widehat{\mu})\) 的）与乘以 \(g\) 的乘法自同态交换，对每个满足下列条件的函数 \(g\in \mathcal{C}_0(\widehat{\mathrm{G}})\)：它属于 \(\mathrm{L}^1 (\mathrm{G},\mu)\) 的傅里叶变换在 \(\mathcal{C}_0(\widehat{\mathrm{G}})\) 中的像（II, 第 223 页, 命题 14）。由于傅里叶变换的像在 \(\mathcal{C}_0(\widehat{\mathrm{G}})\) 中稠密（II, 第 209 页, 命题 5 的推论），态射 \(g\mapsto m_g\) 的连续性蕴含 \(q_{\mathrm{E}}\) 与 \(m_g\) 交换，对每个函数 \(g\in \mathcal{C}_0(\widehat{\mathrm{G}})\)。

于是由 IV 第 188 页的命题 7 的推论，存在函数 \(\varphi \in \mathcal{L}^{\infty}(\widehat{\mathrm{G}},\widehat{\mu})\) 使得 \(q_{\mathrm{E}} = m_{\varphi}\)。由于 \(q_{\mathrm{E}} = q_{\mathrm{E}}^{2}\)，我们有 \(m_{\varphi} = m_{\varphi}^{2} = m_{\varphi^{2}}\)，从而 \(\varphi = \varphi^{2}\)（在 \(\mathrm{L}^{\infty}(\widehat{\mathrm{G}},\widehat{\mu})\) 中）（IV, 第 186 页, 命题 5）；这意味着 \(\varphi\) 在 \(\mathrm{L}^{\infty}(\widehat{\mathrm{G}},\widehat{\mu})\) 中等于 \(\widehat{\mathrm{G}}\) 的某个可测子集 N 的特征函数所在的类。令 \(\mathrm{M} = \widehat{\mathrm{G}} - \mathrm{N}\)。设 \(f \in \mathrm{L}^{2}(\mathrm{G},\mu)\)。我们有：\(f \in \mathrm{E}\) 当且仅当 \(p_{\mathrm{E}}(f) = f\)，当且仅当 \(\varphi \mathcal{F}_{\mathrm{G}}(f) = \mathcal{F}_{\mathrm{G}}(f)\)，而这等价于 \(f \in \mathrm{E}_{\mathrm{M}}\)。

注记. --- 设 \(\chi\) 是 G 的一个特征标。若 G 不是紧的，则 \(\chi\) -同型分量（G 在 L \(^{2}\) (G) 中的正则表示的）

是平凡的。若 G 是紧的，则 G 的正则表示的 \(\chi\) -同型分量的维数为 1，且函数 \(\chi \in \mathrm{L}^{2}(\mathrm{G})\) 是它的一个基。

